\tocsection{\textbf{9/2/26 Lecture 4: }Measures and $\sigma$-Algebras}

Today, we will construct a measure on a $\sigma$-algebra using an outer measure.

\begin{example}{Examples}
    \begin{enumerate}
        \item For any set $X$, let $\mathcal{A} = \P(X)$. Then, we can let $\mu(A) = \#A$, where $\#A$ is the cardinality of $A$. We call $\mu$ the counting measure
        \item Let $X$ and $\mathcal{A} = \P(X)$. Then, fix $x \in X$, and define
        \[ \delta_{x}(A) = \left\{\begin{matrix} 1 & x \in A \\ 0 & \text{otherwise}\end{matrix}\right. \]
        Then, $\delta_{x}(A)$ is a measure, called a point mass.
    \end{enumerate}
\end{example}

\begin{lemma}[\textbf{Proposition 4.1}]
    Let $(X, \cal{A}, \mu)$ be a measure space. Then, the following hold.
    \begin{enumerate}
        \item Let $\{A_n\}_{n = 1}^{N}$ be a finite subset of $\mathcal{A}$ that are pairwise disjoint. Then, we have
        \[ \mu\left(\bigcup_{n = 1}^{N} A_{n}\right) = \sum_{n = 1}^{N} \mu(A_n) \]
        \item If $A, B \in \cal{A}$ and $A \subseteq B$, then $\mu(A) \leq \mu(B)$
        \item If $A_n \in \mathcal{A}$ for each $n \in \N$ not necessarily disjoint, then
        \[ \mu\left(\bigcup_{n = 1}^{\infty} A_n\right) \leq \sum_{n = 1}^{\infty} \mu(A_n) \]
        \item If $A_{1} \subseteq A_2 \subseteq \cdots$ are an increasing sequence of sets in $\cal{A}$, then
        \[ \mu\left(\bigcup_{n = 1}^{\infty} A_n\right) = \lim_{n \to \infty} \mu(A_n) \]
        \item If $A_1 \supseteq A_2 \supseteq \cdots$ is a decreasing sequence in $\cal{A}$ and $\mu(A_1) < \infty$, then
        \[ \mu\left(\bigcap_{n = 1}^{\infty} A_n\right) = \lim_{n \to \infty} \mu(A_n) \]
    \end{enumerate}
\end{lemma}
\begin{proof}{Proof:}
    
    \textbf{1)} Given $A_1, \ldots, A_N$, set $A_m = \emptyset$ for $m > N$. Then, we see that
    \[ \mu\left(\bigcup_{n = 1}^{\infty} A_n\right) = \sum_{n = 1}^{\infty} \mu(A_n) = \sum_{n = 1}^{N} \mu(A_n) \]
    and as $\bigcup_{n = 1}^{\infty} A_n = \bigcup_{n = 1}^{N} A_n$, we have our desired result.\\

    \textbf{2)} Since $B = A \cup (B \backslash A)$, then we see that
    \[ \mu(B) = \mu(A) + \mu(B \backslash A) \]
    as $A$ and $B \backslash A$ are disjoint. Then, as $\mu(B \backslash A) \geq 0$, we have that $\mu(B) \geq \mu(A)$.\\

    \textbf{3)} Define $B_1 = A_1$, $B_2 = A_2 \backslash A_1$ and $B_{3} = A_3 \backslash \left(A_1 \cup A_2\right)$. In general, let $B_{n} = A_n \backslash \left(\bigcup_{i = 1}^{n - 1} A_i\right)$. Now, we see that the $B_{i}$ are disjoint, so
    \[ \mu\left(\bigcup_{n = 1}^{\infty} A_{n}\right) = \mu\left(\bigcup_{n = 1}^{\infty} B_n\right) = \sum_{n = 1}^{\infty} \mu(B_n) \]
    Then, since $\mu(B_n) \leq \mu(A_n)$ as $B_n \subseteq A_n$, we see that
    \[ \mu\left(\bigcup_{n = 1}^{\infty} A_n\right) \leq \sum_{n = 1}^{\infty} \mu(A_n) \]

    \textbf{4)} This is like $5$ so we will do $5$.\\

    \textbf{5)} Let $C_{n} = A_1 \backslash A_n$. By $4$, since $C_{n}$ is increasing, we see that
    \[ \mu\left(\bigcup_{n = 1}^{\infty} C_n\right) = \lim_{n \to \infty} \mu(C_n) = \lim_{n \to \infty} \left( \mu(A_1) - \mu(A_n)\right) = \mu(A_1) - \lim_{n \to \infty}\left(\mu(A_n)\right) \]
    However, we see that
    \[ \bigcup_{n = 1}^{\infty} C_n = A_1 \backslash \bigcap_{n = 1}^{\infty} A_n \]
    so
    \[ \mu\left(\bigcup_{n = 1}^{\infty} C_n\right) = \mu(A_1) - \mu\left(\bigcap_{n = 1}^{\infty} A_n\right) \]
    Thus, we have our desired equality.

    $\qed$
\end{proof}

Note that the condition that $\mu(A_1)$ has finite measure is necessary in $5$. This can be seen by the following counterexample. Take the counting measure on $\R$, and consider the sequence of decreasing sets
\[ A_n = [n, \infty) \]
Then, we see that $\mu(A_n) = \infty$ for each $n$, but $\bigcap_{n} A_n = \emptyset$.

\begin{definition}[\textbf{Definition 4.2:} $\sigma$-algebra Generated by a Collection]
    Let $\cal{C}$ be a collection of subsets of a set $X$. Then, the \emph{$\sigma$-algebra generated by $\cal{C}$} is 
    \[ \sigma(\cal{C}) = \bigcap \{\cal{A} \col \cal{A} \text{ is a } \sigma \text{-algebra containing } \cal{C}\} \]
    Since the intersection of $\sigma$-algebras is a $\sigma$-algebra, $\sigma(\cal{C})$ is indeed a $\sigma$-algebra.
\end{definition}
In a sense, we say that $\sigma(\cal{C})$ is the smallest $\sigma$-algebra containing $\cal{C}$ (since any $\sigma$-algebra containing $\cal{C}$ must contain $\sigma(\cal{C})$). 

\begin{definition}[\textbf{Definition 4.3:} Borel $\sigma$-algebra]
    Let $X$ be a topological space, and let $G$ be it's set of open sets. Then, the \emph{Borel $\sigma$-algebra} is
    \[ \cal{B} := \sigma(G) \]
    We call $A \in \cal{B}$ a \emph{Borel set}.
\end{definition}

\begin{lemma}[\textbf{Proposition 4.4}]
    Consider $\R$ with the normal topology, and let $\cal{B}$ be the Borel $\sigma$-algebra. Then, $\cal{B}$ is generated by each of the following collections of subsets.
    \begin{enumerate}
        \item $\cal{C}_1 = \{(a, b) \col a, b \in \R\}$
        \item $\cal{C}_2 = \{[a, b] \col a, b \in \R\}$
        \item $\cal{C}_3 = \{(a, b] \col a, b \in \R\}$
        \item $\cal{C}_4 = \{(a, \infty) \col a \in \R\}$
    \end{enumerate}
\end{lemma}
\begin{proof}{\textbf{Proof:}}
    We won't do all of them, but we will show a few. The idea is to show that we can either generate the open sets with this collection, or the other way around.\\

    \textbf{1)} First, any open set can be written as $\cup_{n = 1}^{\infty} I_n$ where $I_n = (a, b), (-\infty, a)$ or $(a, \infty)$. Now, we can write
    \[ (a, \infty) = \cup_{n} (a, a + n) \]
    and similarly, 
    \[ (-\infty, a) = \cup_{n} (a - n, a) \]
    Thus, $(a, b), (-\infty, a)$ and $(a, \infty)$ are in $\sigma(\cal{C}_1)$. Thus, we see that $\cal{B} \subseteq \sigma(\cal{C}_1)$. Since $C_1 \subseteq \cal{B}$, we equality.\\
    
    \textbf{2)} The goal is to show that $\sigma(\cal{C}_2) \subseteq \cal{B}$ and that $\cal{B} \subseteq \sigma(\cal{C}_2)$. Since
    \[ [a, b] = \bigcap_{n} (a - 1/n, b + 1/n) \]
    we see that $\cal{C}_2 \subseteq \sigma({\cal{C}_1}) = \cal{B}$. Thus, $\sigma({\cal{C}_2}) \subseteq \cal{B}$. Then, as
    \[ (a, b) = \bigcup_{n} [a + 1/n, b - 1/n] \]
    we see that $\cal{C}_1 \subseteq \sigma(\cal{C}_2)$, so $\cal{B} = \sigma(\cal{C}_1) \subseteq \sigma(\cal{C}_2)$. Thus, we have equality.\\

    \textbf{3)} and \textbf{4)} are very similar, so we leave the proof to the reader.

    $\qed$
\end{proof}

\begin{definition}[\textbf{Definition 4.5:} $\mu^{*}$-Measurability]
    For an outer measure $\mu^{*}$ on a set $X$, we sey that $A \subseteq X$ is \emph{$\mu^{*}$-measurable} if, for any $E \subseteq X$, then
    \[ \mu^{*}(E) = \mu^{*}(E \cap A) + \mu^{*}(E \cap A^{c}) \]
\end{definition}
